\section{Architecture}
\label{sec:architecture}

\subsection{Name and structural principle}

\textbf{ALUCLU} expands to \textbf{Adaptive Layered Unified Context with
Learned Updates}. ``Layered'' denotes scheduled placement across depth and the
separation of physical memory regimes. ``Unified context'' denotes a single
token representation assembled from all active lanes. ``Learned updates''
refers both to recurrent write dynamics and token-wise fusion. The name is
written with ASCII \texttt{C} in code, package names, and URLs.

The unit of composition is \texttt{AlucluBlock}. Given input \(x_t^{(\ell)}\),
the block first computes
\begin{equation}
  h_t^{(\ell)}=\RMSNorm_{\mathrm{mix}}\!\left(x_t^{(\ell)}\right).
\end{equation}
Every active memory lane reads \(h_t^{(\ell)}\) in parallel. A lane never
consumes another lane's output as its input. This keeps lane semantics
separable and lets each state transition be audited against the same token.

\subsection{Layer scheduling}

Every one of the \(L\) layers contains a Gated Delta lane. The other lanes are
periodic:
\begin{align}
  N_{\mathrm{local}}
  &=\left\lfloor\frac{L}{e_{\mathrm{local}}}\right\rfloor,\\
  N_{\mathrm{epi}}
  &=\left\lfloor\frac{L}{e_{\mathrm{epi}}}\right\rfloor,\\
  N_{\mathrm{exact}}
  &=\left\lfloor\frac{L}{e_{\mathrm{exact}}}\right\rfloor,
  \label{eq:branch-counts}
\end{align}
when the corresponding configuration exists. Here each \(e\) is the
``every'' interval. This schedule permits a recurrent backbone at every depth
while concentrating higher-state or higher-latency lanes at selected blocks.

\subsection{Block data flow}

\Cref{fig:block} shows the implemented data flow. The residual-surprise cache
receives its scalar insertion priority from the Gated Delta write residual,
but its query, key, and value projections still read \(h_t^{(\ell)}\).

\begin{center}
\resizebox{\textwidth}{!}{%
\begin{tikzpicture}[
  node distance=6mm and 8mm,
  every node/.style={font=\small},
  box/.style={draw=AlucluNavy,rounded corners=2pt,thick,fill=white,
    align=center,minimum height=9mm},
  lane/.style={box,text width=35mm},
  norm/.style={box,fill=AlucluLight,text width=25mm},
  sum/.style={circle,draw=AlucluNavy,thick,fill=AlucluLight,
    minimum size=8mm,inner sep=0pt},
  arrow/.style={-{Latex[length=2.2mm]},thick,color=AlucluNavy},
  signal/.style={-{Latex[length=2mm]},dashed,color=AlucluGold,thick}
]
\node[box,text width=28mm] (input) {\(x_t^{(\ell)}\)};
\node[norm,right=of input] (mixnorm) {Mixer\\RMSNorm};

\node[lane,above right=10mm and 17mm of mixnorm] (gdn)
  {Gated Delta Rule-2\\every layer};
\node[lane,below=of gdn] (local)
  {Exact local attention\\last \(W\) tokens};
\node[lane,below=of local] (cache)
  {Residual-surprise cache\\top \(C_x\) projected KV};
\node[lane,below=of cache] (epi)
  {Episodic hierarchy\\current + live + archive};

\node[norm,right=12mm of gdn] (ng) {Branch norm\\and scale};
\node[norm,right=12mm of local] (nl) {Branch norm\\and scale};
\node[norm,right=12mm of cache] (nc) {Branch norm\\and scale};
\node[norm,right=12mm of epi] (ne) {Branch norm\\and scale};

\node[box,fill=AlucluLight,text width=29mm,right=13mm of nl,
  yshift=-9mm] (fusion) {Token-wise softmax\\fusion};
\node[sum,right=10mm of fusion] (add1) {\(+\)};
\node[norm,right=8mm of add1] (ffnnorm) {FFN\\RMSNorm};
\node[box,text width=22mm,right=8mm of ffnnorm] (swiglu) {SwiGLU};
\node[sum,right=8mm of swiglu] (add2) {\(+\)};
\node[box,text width=28mm,right=8mm of add2] (output)
  {\(x_t^{(\ell+1)}\)};

\draw[arrow] (input) -- (mixnorm);
\draw[arrow] (mixnorm) |- (gdn);
\draw[arrow] (mixnorm) |- (local);
\draw[arrow] (mixnorm) |- (cache);
\draw[arrow] (mixnorm) |- (epi);
\draw[signal] (gdn) -- node[above,text=AlucluGold]
  {\(\xi_t\)} (cache);

\draw[arrow] (gdn) -- (ng);
\draw[arrow] (local) -- (nl);
\draw[arrow] (cache) -- (nc);
\draw[arrow] (epi) -- (ne);
\draw[arrow] (ng) -| (fusion);
\draw[arrow] (nl) -- (fusion);
\draw[arrow] (nc) -- (fusion);
\draw[arrow] (ne) -| (fusion);
\draw[arrow] (mixnorm) -- ++(0,-36mm) -| node[pos=0.82,below]
  {fusion logits} (fusion);
\draw[arrow] (fusion) -- (add1);
\draw[arrow] (add1) -- (ffnnorm);
\draw[arrow] (ffnnorm) -- (swiglu);
\draw[arrow] (swiglu) -- (add2);
\draw[arrow] (add2) -- (output);
\draw[arrow] (input) to[out=90,in=90,looseness=1.45] (add1);
\draw[arrow] (add1) to[out=-90,in=-90,looseness=1.35] (add2);
\end{tikzpicture}
}
\captionof{figure}{\ALUCLU block. Solid arrows carry vectors. The dashed arrow carries
the scalar Gated Delta write-surprise signal used by the exact-cache
replacement policy. Every memory lane otherwise reads the same normalized
token.}
\label{fig:block}
\end{center}

\subsection{Streaming causal stems}

The Gated Delta and episodic lanes each begin with an independent causal
depthwise convolution of kernel width \(\kappa\). For an input chunk,
\begin{align}
  c &= [h_{\mathrm{history}};x_{1:T}],\\
  \widetilde h_{1:T} &= \operatorname{DWConv1D}(c),\\
  h_{\mathrm{history}}' &= c_{-(\kappa-1):}.
  \label{eq:stream-conv}
\end{align}
The persistent history has \(B(\kappa-1)d\) elements. Kernels are initialized
to identity by setting all weights to zero except the final causal tap, which
is one, with zero bias. Thus adding the stem does not perturb the initial
token path while allowing training to learn short causal filtering.

\subsection{Token-wise branch fusion}
\label{sec:fusion}

Let \(b_{t,j}\) be active branch \(j\)'s output. Each branch has a separate
normalizer and learned scalar:
\begin{equation}
  \widetilde b_{t,j}
  =s_j\,\RMSNorm_j(b_{t,j}).
\end{equation}
For more than one branch, the mixture is
\begin{align}
  \alpha_t &= \softmax(W_fh_t+b_f),\\
  f_t &= \sum_j\alpha_{t,j}\widetilde b_{t,j},
  \qquad \alpha_{t,j}>0,\quad\sum_j\alpha_{t,j}=1.
  \label{eq:fusion}
\end{align}
The coefficients form a convex combination of the \emph{scaled normalized}
branches. Because \(s_j\) is unconstrained, \(f_t\) is not necessarily in the
convex hull of the raw \(b_{t,j}\). The projection is initialized to zero, so
the initial coefficients are uniform; all branch scales begin at one.

The residual and feed-forward path is
\begin{align}
  x_t' &= x_t+\operatorname{Dropout}(f_t),\\
  (g_t,v_t) &=
  \operatorname{split}\!\left(W_{\mathrm{in}}\RMSNorm(x_t')\right),\\
  u_t &= W_{\mathrm{out}}\!\left(\SiLU(g_t)\odot v_t\right),\\
  x_t^{\mathrm{out}} &= x_t'+\operatorname{Dropout}(u_t).
  \label{eq:swiglu}
\end{align}

\subsection{State hierarchy}

The public model state is a tuple of per-block states. A block state contains
the always-present Gated Delta state and optional local, episodic, and
exact-cache states. The state is returned explicitly by \texttt{forward} and
\texttt{step}; it is never hidden inside the module or reset implicitly.

\begin{center}
\begin{tikzpicture}[
  level distance=11mm,
  sibling distance=37mm,
  every node/.style={font=\small,draw=AlucluNavy,rounded corners,
    fill=AlucluLight,align=center,inner sep=3pt},
  edge from parent/.style={draw=AlucluNavy,-{Latex[length=1.5mm]}}
]
\node {LanguageModelState}
  child {node {BlockState \(\times L\)}
    child {node {GatedDeltaState\\history, fast weight}}
    child {node {LocalAttentionState\\rolling K/V}}
    child {node {ExactCacheState\\K/V, priority, valid}}
    child {node {EpisodicState\\current, live, archive}}};
\end{tikzpicture}
\captionof{figure}{Logical persistent-state ownership. Optional branches appear only in
scheduled layers.}
\label{fig:state-tree}
\end{center}

\subsection{Sequence, chunk, and step paths}

The same token recurrence backs three interfaces:

\begin{description}
  \item[Full sequence.] Consume a complete \(B\times T\) tensor from an
  optional incoming state and return outputs plus final state.
  \item[Chunk.] Invoke the full-sequence path repeatedly while carrying the
  returned state. This is useful for truncated backpropagation or streaming
  prefill.
  \item[Step.] Consume one token per batch item and return one-token logits and
  state.
\end{description}

In evaluation mode, the tests require these partitions to produce equivalent
outputs and terminal state within the declared numerical tolerance. Training
mode with dropout is not partition invariant because different partitions
consume random numbers in a different order.
